\chapter{I\textsuperscript{2}C Example}
\label{chap:fpgasteel-i2c-example}

\section{Prerequisites}
\label{sec:fpgasteel-i2c-prerequisites}

Same toolchain, \term{SD card} and \term{Eclipse}/\term{OpenOCD} debug setup as Hello World (Chapter~\ref{chap:fpgasteel-hello-world}).\footnote{TODO: footnote text was lost when the source was copy-pasted --- please supply it.} In addition, the \term{OV7670} needs the custom camera adapter board (line buffer between the camera module and the \term{FPGA}) plugged into the \term{GPIO\_0} header on the \term{DE1-SoC}; schematics, bill of materials and photographs are in Appendix~\ref{app:camera-adapter-board}; the \term{PCB} files are under |hw/camera_adapter|.

\begin{warningbox}
    \textbf{Hardware bug:} do not leave the camera in \term{RESET} or \term{POWER-DOWN} for extended periods.\footnote{TODO: footnote text was lost when the source was copy-pasted --- please supply it.}
\end{warningbox}

\section{About this project}
\label{sec:fpgasteel-i2c-about}

|i2c_example| (under |sw/i2c_example|) checks that the \term{HPS} can communicate with the \term{OV7670} camera module over its control bus (\term{SCCB}, an \term{I\textsuperscript{2}C}-like protocol) before any actual camera driver is written. It:

\begin{itemize}
    \item brings up the interrupt controller and \term{UART0}, as in |hello_word|;
    \item releases the camera from reset;
    \item reads the \term{OV7670}'s \term{PID} register over \term{SCCB} and prints it;
    \item spins forever.
\end{itemize}

On a working setup, the serial console prints |OV7670 PID = 0x76|, the \term{OV7670}'s fixed product \term{ID} (datasheet, register 0x0A).

\section{FPGA Hardware Involved}
\label{sec:fpgasteel-i2c-hardware}

Software communicates with the camera over \term{SCCB} through two modules placed in \term{Platform Designer} (|hw/quartus/soc_system.qsys|), each under its own instance name:

\begin{itemize}
    \item \textbf{\texttt{ov7670\_ctrl\_interface}:} an instance of \term{Altera}/\term{Intel}'s |altera_avalon_i2c| \term{IP}, used as the \term{I\textsuperscript{2}C}/\term{SCCB} master.
    \item \textbf{\texttt{ov7670\_PCLK\_reset\_controller}:} a custom \term{IP} (|pclk_reset_controller|, inherited unchanged from \sysname{FPGAlix} and renamed here, not part of \term{Altera}/\term{Intel}'s \term{IP} catalog) that holds the camera in reset. While in reset, the camera doesn't drive \term{PCLK} either, so the capture logic downstream (|ov7670_video_interface|) and the dual-clock \term{FIFO} right after it on the \term{Avalon-ST} path (|ov_7670_video_dc_fifo|) have no clock to run on. The reset must be released before attempting \term{SCCB} communication or image acquisition, or the camera won't respond. Renamed from \sysname{FPGAlix}'s |pclk_reset_controller_0|, see Paragraph~\ref{subsec:fpgasteel-renamed-instances}. The two-reset scheme it implements is covered in full in PART~I -- \sysname{FPGAlix}, Section~\ref{sec:fpgalix-pclk-reset}.
\end{itemize}

\section{Code structure}
\label{sec:fpgasteel-i2c-code-structure}

\subsection{AvalonI2C}
\label{subsec:fpgasteel-i2c-avaloni2c}

|AvalonI2C| class (|src/AvalonI2C.{h,cpp}|) is a polling-mode driver for the \term{Avalon I\textsuperscript{2}C} \term{IP} (|altera_avalon_i2c|), master mode, 7-bit addressing only. Unlike |System|/|Serial| it is not a singleton: a board can have several \term{Avalon I\textsuperscript{2}C} cores, so each instance just wraps a register base address passed to its constructor. Its public methods are:

\begin{itemize}
    \item \textbf{\texttt{init(speed)}:} |speed| selects the bus speed mode, |Speed::Standard| (100 \term{kHz}) or |Speed::Fast| (400 \term{kHz}). Disables the core (|CTRL|), masks interrupts (|ISER|, this driver polls, it doesn't use them), clears pending status bits (|ISR|), programs the |SCL_HIGH|/|SCL_LOW|/|SDA_HOLD| bus-timing registers for the requested speed, then reconfigures |CTRL| with the bus speed and \term{FIFO} thresholds and enables the core.
    \item \textbf{\texttt{deInit()}:} waits for the core to go idle, then disables it.
    \item \textbf{\texttt{write(devAddress, regAddress, regVal)}:} |devAddress| is the 7-bit \term{I\textsuperscript{2}C} slave address, |regAddress| the register to write within that slave, |regVal| the value to write. Runs the common \term{I\textsuperscript{2}C} ``device address, register address, value'' sequence (\term{START}, dev address+W, register address, value, \term{STOP}).
    \item \textbf{\texttt{read(devAddress, regAddress)}:} |devAddress| is the 7-bit \term{I\textsuperscript{2}C} slave address, |regAddress| the register to read; returns the value read. Writes the register address on the \term{I\textsuperscript{2}C} bus, then a repeated \term{START} into a read (\term{START}, dev address+W, register address, \term{RESTART}, dev address+R, value, \term{STOP}): standard \term{I\textsuperscript{2}C} register read.
    \item \textbf{\texttt{isReady()}:} whether |init()| has run and |deInit()| hasn't since.
\end{itemize}

Internally, both |write()| and |read()| throw |ExceptionI2C104| on a \term{NACK}, an arbitration loss, or if the core doesn't respond within |CMD_TIMEOUT_MS| (10 \term{ms}): this is what turns a broken link into an error message instead of a hang.

\subsection{AvalonSCCB}
\label{subsec:fpgasteel-i2c-avalonsccb}

|AvalonSCCB| class (|src/AvalonSCCB.{h,cpp}|) derives from |AvalonI2C| and only overrides |read()|. \term{SCCB} differs from \term{I\textsuperscript{2}C} in exactly this: instead of one transaction with a repeated \term{START}, \term{SCCB} requires two separate \term{START}\ldots\term{STOP} transactions: (1) \term{START}, dev address+W, register address, \term{STOP}; (2) \term{START}, dev address+R, value, \term{STOP}. |write()| is unchanged, since \term{SCCB} and \term{I\textsuperscript{2}C} writes are identical, so it's inherited as-is from |AvalonI2C|.

\subsection{PclkResetController}
\label{subsec:fpgasteel-i2c-pclkresetcontroller}

|PclkResetController| class (|src/PclkResetController.{h,cpp}|) is a singleton (there's only one on the board) wrapping the |pclk_reset_controller| core's single control register\footnote{TODO: footnote text was lost when the source was copy-pasted --- please supply it.}:

\begin{itemize}
    \item \textbf{\texttt{assertReset()}:} holds the |ov7670_pclk| domain and the camera in reset (|sw_release| = 0).
    \item \textbf{\texttt{deassertReset()}:} releases it (|sw_release| = 1).
\end{itemize}

\subsection{ExceptionI2C}
\label{subsec:fpgasteel-i2c-exceptioni2c}

|ExceptionI2C| class (|src/Exceptions/ExceptionI2C.h|) follows the same pattern as |ExceptionSystem|/|ExceptionSerial| (Chapter~\ref{chap:fpgasteel-hello-world}, Section~\ref{sec:fpgasteel-hello-code-structure}), but carries an |ErrorType| instead of an |ALT_STATUS_CODE|: |Nack|, |ArbitrationLost|, |Timeout|, or |NotInitialized| (a call made before |init()|).

\section{How it works}
\label{sec:fpgasteel-i2c-how-it-works}

The following snippet is taken from |main.cpp|:

\begin{lstlisting}[caption={The main() body of i2c\_example.}, label={lst:fpgasteel-i2c-main}]
System::getInstance().init();
Serial::getInstance().init();
PclkResetController::getInstance().deassertReset();
// release the camera from reset
AvalonSCCB sccb(0xFF200180);
// ov7670_ctrl_interface
sccb.init();
uint8_t pid = sccb.read(0x21, 0x0A);
// OV7670 SCCB address, PID register
serial->printf("OV7670 PID = 0x%02X\r\n", pid);
\end{lstlisting}

|ExceptionSystem|, |ExceptionSerial|, |ExceptionRingBuffer| and |ExceptionI2C| are all caught around this block and reported over \term{Serial}, exactly as in |hello_word|.

\section{References}

\begin{description}
    \item[EP-IPUG] Chapter 15, ``Intel FPGA Avalon \term{I\textsuperscript{2}C} (Host) Core'': register map (Section 15.5) and programming model (Section 15.7.6, Figure 54), followed by |AvalonI2C::init()|.
    \item[OV7670-DS] \term{PID} register (0x0A, fixed value 0x76) used as the check in Section~\ref{sec:fpgasteel-i2c-how-it-works}.
\end{description}

